Las cifras de esta sección (500~ms, 1~min, 99,5\,\%) son \emph{objetivos de diseño}; se
comprobarán con pruebas en PE-U1.

\paragraph{Recorrido de un mensaje directo.}
Al abrir la aplicación, el cliente (CLI) establece un canal WebSocket seguro (WSS sobre TCP)
con el balanceador, que lo asigna a una de las instancias de MSG (S4). El canal queda
abierto en ambos sentidos. Un mensaje pasa por tres tramos con sincronía distinta:
\begin{enumerate}[leftmargin=*]
  \item \textbf{Envío y ACK de guardado (síncrono para el remitente).} El cliente genera un
  \texttt{message\_id} único y envía el mensaje; MSG lo guarda en PostgreSQL (S6) y solo
  entonces responde al remitente con el \textbf{ACK de guardado}. El remitente espera ese
  ACK y, si no llega, reenvía con el mismo \texttt{message\_id}.
  \item \textbf{Entrega al destinatario (asíncrona).} Después de guardar, MSG empuja el
  mensaje por el canal WSS del destinatario si está conectado. El ACK al remitente no
  espera a esta entrega.
  \item \textbf{ACK de entrega (asíncrono).} El cliente del destinatario confirma la
  recepción con el \texttt{sequence\_no} recibido; MSG registra ese valor como el último
  confirmado para ese destinatario en esa conversación.
\end{enumerate}
Por tanto, el \textbf{ACK al remitente significa ``guardado en PostgreSQL'', no ``entregado''}.
La entrega se confirma aparte con el ACK de entrega del destinatario.

\paragraph{Presencia en línea.}
Mientras el canal está abierto, el cliente envía periódicamente un \emph{heartbeat} por el
mismo canal WSS (S5); no es una petición que bloquee al cliente, por lo que se trata como
asíncrono y periódico. Con cada uno, MSG renueva en Redis la clave de presencia del usuario,
que tiene un tiempo de expiración (TTL) (S7). Si los \emph{heartbeats} dejan de llegar, la
clave expira y el usuario se muestra como desconectado, sin necesidad de un aviso explícito.

\paragraph{Orden: asignación de \texttt{sequence\_no} entre varias instancias de MSG.}
Como hay varias instancias, el número de secuencia no puede asignarse en memoria. Se asigna
en PostgreSQL, que es el punto común: cada conversación tiene un contador y, dentro de una
misma transacción, MSG incrementa ese contador y guarda el mensaje con el valor obtenido. El
bloqueo de la fila del contador hace que dos instancias que reciben mensajes de la misma
conversación al mismo tiempo se ordenen una detrás de otra y no puedan obtener el mismo
número. La tabla de mensajes impone que el par (\texttt{conversation\_id},
\texttt{sequence\_no}) y el \texttt{message\_id} sean únicos.

\paragraph{Deduplicación y cursor de recuperación.}
Antes de asignar un número, MSG busca el \texttt{message\_id}: si ya existe (por ejemplo, un
reintento tras perder el ACK), devuelve el \texttt{sequence\_no} ya asignado y no crea otro
mensaje ni avanza el contador. Para recuperarse, el cliente conserva por conversación el
último \texttt{sequence\_no} confirmado y, al reconectarse (F1 o F3), lo envía como cursor;
MSG responde con los mensajes de esa conversación cuyo \texttt{sequence\_no} es mayor,
en orden. Si el cliente detecta un salto en la numeración, pide el rango que falta.

\paragraph{Denuncia y ocultamiento.}
Cuando un usuario denuncia un aviso, el cliente envía a MOD el \texttt{id\_aviso}, el
\texttt{id\_usuario} y el motivo (S8), y MOD confirma la recepción. MOD cuenta únicamente las
denuncias válidas de usuarios distintos. Al llegar a tres, solicita a PUB ocultar el aviso
(S9); PUB cambia su estado a oculto, que es la fuente de verdad, y publica en Kafka el evento
\texttt{PublicaciónOcultada} (A4), que MUR consume para retirar o invalidar el aviso en los
muros.

\paragraph{Mecanismos y justificación.}
\begin{itemize}[leftmargin=*]
  \item \textbf{WSS sobre TCP, bidireccional y persistente (S4, S5).} Un canal abierto
  permite empujar el mensaje al destinatario apenas se guarda, sin abrir una conexión por
  cada envío, y sirve también para el \emph{heartbeat}. TCP garantiza entrega ordenada y
  confiable; WSS cifra el tráfico.
  \item \textbf{PostgreSQL síncrono antes del ACK (S6).} Un mensaje no se confirma al
  remitente hasta estar guardado, para no confirmar algo que luego pueda perderse.
  \item \textbf{Redis con TTL (S7).} La presencia es un dato temporal que cambia
  constantemente y no necesita sobrevivir a un reinicio, por lo que se guarda en memoria con
  expiración.
  \item \textbf{REST y gRPC en la moderación (S8, S9) y Kafka en la invalidación (A4).} La
  denuncia necesita confirmación inmediata; la orden de ocultar es una llamada directa entre
  servicios internos; y el retiro del aviso en los muros puede completarse de forma
  asíncrona.
\end{itemize}

\paragraph{R2: objetivo de 500~ms y orden.}
El canal WSS ya está abierto, por lo que la entrega a un destinatario conectado no depende
de establecer conexiones nuevas y la tarea más costosa de la ruta es una escritura en
PostgreSQL. El \texttt{sequence\_no} asignado en PostgreSQL mantiene el orden por
conversación aunque haya varias instancias, y el \texttt{message\_id} único evita duplicados.
Se plantea como objetivo y se medirá en PE-U1.

\paragraph{R5: objetivo de ocultar en menos de un minuto.}
La solicitud a PUB (S9) es síncrona y cambia el estado de inmediato. El evento A4 es
asíncrono, así que por sí solo no garantiza el límite: si el evento aún no se consumió, MUR
podría seguir mostrando su copia del aviso. Para cerrar ese hueco, MUR no se fía solo del
evento: al armar el muro para un usuario, consulta a PUB en lote el estado de los avisos que
va a mostrar, con una caché de pocos segundos, y descarta los ocultos. Así, tras el cambio de
estado en PUB, el aviso deja de verse en cuanto expira esa caché, y el evento A4 se encarga
después de limpiar los muros ya armados. Se plantea como objetivo y se medirá en PE-U1.

\paragraph{R6: objetivo de disponibilidad de 99,5\,\% en la mensajería.}
Se aspira a esa meta con varias instancias de MSG detrás del balanceador, \emph{health
checks} que retiran las instancias que no responden y persistencia de los mensajes en
PostgreSQL. El muro está desacoplado mediante Kafka, de modo que su degradación no bloquea
la mensajería.

\paragraph{F1: el destinatario está desconectado.}
El mensaje ya se guardó en PostgreSQL antes del ACK de guardado, por lo que queda pendiente
aunque el destinatario no esté en línea. Al reconectarse, el cliente envía su cursor (último
\texttt{sequence\_no} confirmado por conversación) y MSG le entrega los mensajes posteriores.

\paragraph{F3: MSG se reinicia con conversaciones abiertas.}
Los clientes se reconectan al balanceador, que los asigna a una instancia disponible, y
recuperan desde su cursor, sin depender de memoria volátil. Los mensajes que el remitente no
vio confirmados se reenvían con el mismo \texttt{message\_id}, y la deduplicación evita que
se guarden dos veces.
